\chapter{Conclusiones}\label{conclusion}

\azul{Este capítulo finaliza el documento formalizando el cierre de la investigación o desarrollo ingenieril. Su propósito no es resumir el trabajo paso a paso, sino sintetizar los hallazgos definitivos, el nivel de éxito del artefacto o sistema construido y las implicancias prácticas del estudio. La literatura metodológica establece que en esta sección está estrictamente prohibido introducir nuevos datos, ecuaciones, citas bibliográficas o conceptos teóricos que no hayan sido analizados en los capítulos precedentes \cite[p.~360]{hernandez2014metodologia}.}

\section{Principales Aportes de la Ingeniería}

\azul{Se requiere detallar las contribuciones técnicas, sociales, económicas o de infraestructura más relevantes que la solución aporta al entorno de aplicación. Se debe enfatizar el valor tangible o intangible del producto entregable, independientemente de si es una obra civil, un plan de optimización industrial, o una plataforma tecnológica.}

\moradoc{Ejemplo: El principal aporte de este proyecto radica en la entrega de un modelo matemático validado que reduce la incertidumbre en el cálculo de cargas sísmicas para estructuras de contención. Adicionalmente, el diseño propuesto optimiza el uso de materiales en un 12\%, generando un impacto económico directo y ambientalmente sostenible para la agroindustria de la Región de Los Lagos.}

\section{Contraste de Resultados y Cumplimiento de Objetivos}

\azul{Esta sección exige una evaluación crítica, reflexiva y directa del grado de cumplimiento de los objetivos específicos planteados en el Capítulo \ref{metodologia}. El/la estudiante debe contrastar explícitamente los resultados obtenidos en las fases de diseño, implementación y validación contra cada meta inicial. Si un objetivo se cumplió solo de forma parcial, se debe argumentar técnicamente el motivo, demostrando madurez profesional.}

\moradoc{Ejemplo: Respecto al Objetivo Específico \ref{oe1} (Diagnosticar las condiciones de la red hídrica), este se logró a cabalidad mediante el levantamiento en terreno de tres nodos críticos. En cuanto al Objetivo Específico \ref{oe4} (Operar y validar el sistema), se reporta un cumplimiento parcial: si bien las pruebas de laboratorio confirmaron la precisión de los instrumentos, las condiciones de contingencia climática impidieron la ejecución del 100\% de las pruebas de estrés en terreno estipuladas originalmente.}

\section{Limitaciones y Trabajos Futuros}

\azul{El apartado final reconoce con honestidad académica las restricciones de diseño, de presupuesto, tecnológicas o temporales que acotaron el alcance del proyecto (limitaciones). A partir de esto, se proponen nuevas líneas de acción para futuras investigaciones, optimizaciones o rediseños. Esta sección sirve como punto de partida para que futuras y futuros tesistas del Departamento de Ciencias de la Ingeniería puedan continuar expandiendo el desarrollo.}

\moradoc{Ejemplo: Entre las limitaciones del diseño actual se identifica que el marco de sensores requiere alimentación eléctrica ininterrumpida, lo que dificulta su despliegue en zonas totalmente aisladas. Como trabajo futuro, se propone a las próximas cohortes de ingeniería el diseño de un subsistema de alimentación basado en paneles fotovoltaicos y el desarrollo de un protocolo de transmisión satelital para zonas sin cobertura de red tradicional.}